% =====================================================================
%  Harald Høffding, MINDRE ARBEJDER. TREDIE RÆKKE.
%  København: Gyldendal, 1913.  234 pp.  (Dansk Bogfortegnelse 1909-14, p. 147: "Mindre Arbejder. 3. Række.
%  (234 S.). 13. Gyldendal. 4,00.")
%
%  STATUS: IN PROGRESS (2026-10-25).  Only the INDHOLD is transcribed so far, as a placeholder until the book can be
%  transcribed whole.  No scan has been found (NB.no lists the volume, Norway-only and not viewable; not on the Internet
%  Archive; HathiTrust and Google Books not yet searched).  Not to be marked complete.
%
%  SOURCE OF THE INDHOLD.  Hans Halvorson's own copy, photographed 2026-10-25:
%  ~/bibliotek/Høffding, Harald/1913-mindre-arbejder-3-indhold.jpg.  Read from the photograph in the main session; the
%  page is clean and every entry, year and page number is legible.  The Indhold is an unnumbered page; its position in
%  the front matter is not yet known.  Set as printed: numbers, titles, the year of each piece in parentheses, page
%  numbers; the dot leaders as \dotfill; a short rule under the heading and one after the list.
% =====================================================================
\documentclass[12pt,a4paper,openany]{book}
\usepackage[utf8]{inputenc}
\usepackage[T1]{fontenc}
\usepackage{libertinus}
\usepackage{libertinust1math}
\usepackage[danish]{babel}
\usepackage[protrusion=true,expansion=true]{microtype}
\usepackage{setspace}
\usepackage[margin=1.3in]{geometry}
\usepackage{fancyhdr}
\usepackage{hyperref}
\hypersetup{pdftitle={Mindre Arbejder. Tredie Række}, pdfauthor={Harald Høffding}, hidelinks}
\pagestyle{fancy}
\fancyhf{}
\fancyhead[LE,RO]{\thepage}
\fancyhead[C]{\textit{Høffding: Mindre Arbejder. Tredie Række (1913)}}
\setlength{\headheight}{15pt}
\setstretch{1.3}
\usepackage{marginnote}
\newcommand{\opage}[1]{\marginnote{\footnotesize\textit{[#1]}}}
\newcommand{\ix}[3]{\par\noindent\makebox[2em][r]{#1}\hspace{0.6em}#2\dotfill #3}

\begin{document}

\frontmatter
\chapter*{INDHOLD}
\begin{center}\rule{2cm}{0.4pt}\end{center}
\noindent\hfill Side
\ix{1.}{Aandelig Kultur (1912)}{1}
\ix{2.}{En filosofisk Bekendelse (1904)}{18}
\ix{3.}{Begrebet Villie (1905; 1912)}{28}
\ix{4.}{Human Etik (1901)}{54}
\ix{5.}{Sociologi og Etik (1904)}{68}
\ix{6.}{En Sygehistorie (1901)}{78}
\ix{7.}{Platons Symposion (1905)}{83}
\ix{8.}{Demokritos og Platon (1909)}{94}
\ix{9.}{Epiktet og Kristendommen (1912)}{102}
\ix{10.}{Marcus Aurelius (1912)}{107}
\ix{11.}{Rousseau og det nittende Aarhundrede (1912)}{120}
\ix{12.}{Rousseau og Religionen (1912)}{142}
\ix{13.}{Goethe og Filosofien (1907)}{162}
\ix{14.}{Ved Schillerfesten (1909)}{176}
\ix{15.}{Napoleon og Filosofien (1906)}{184}
\ix{16.}{Kierkegaard og Nietzsche (1911)}{189}
\ix{17.}{Ved Mindefesten for Krøyer (1909)}{196}
\ix{18.}{Darwin og Filosofien (1909)}{202}
\par\bigskip
\begin{center}\rule{4cm}{0.4pt}\end{center}

\mainmatter
% [the text, pp. 1-234, to be transcribed]

\end{document}
